% Created 2021-07-13 ter 18:00
% Intended LaTeX compiler: pdflatex
\documentclass[presentation]{beamer}
\usepackage[utf8]{inputenc}
\usepackage[T1]{fontenc}
\usepackage{graphicx}
\usepackage{grffile}
\usepackage{longtable}
\usepackage{wrapfig}
\usepackage{rotating}
\usepackage[normalem]{ulem}
\usepackage{amsmath}
\usepackage{textcomp}
\usepackage{amssymb}
\usepackage{capt-of}
\usepackage{hyperref}
\usepackage{graphicx, hyperref, url}
\mode<beamer>{\usetheme{Madrid}}
\usetheme{default}
\author{Rafael Castro - rafaelcgs10.github.io}
\date{14/07/2021}
\title{BBL sobre Nix}
\hypersetup{
 pdfauthor={Rafael Castro - rafaelcgs10.github.io},
 pdftitle={BBL sobre Nix},
 pdfkeywords={},
 pdfsubject={},
 pdfcreator={Emacs 27.2 (Org mode 9.4.6)}, 
 pdflang={Portuguese}}
\begin{document}

\maketitle

\section{Introdução}
\label{sec:org31fa441}
\begin{frame}[label={sec:orgea77b0d}]{O que é cada coisa?}
\begin{itemize}
\item Nix: Um gerenciador de pacotes \alert{puramente funcional}
\begin{itemize}
\item Nix Expression Language: Uma linguagem pura, funcional e lazy
\item NixOS: Uma distro Linux feita envolta do Nix
\item Home Manager: Um Manager (gerenciador) para espaço de usuário envolta do Nix
\end{itemize}
\end{itemize}
\end{frame}

\section{Show of:}
\label{sec:org5293d55}
\begin{frame}[label={sec:org0403419}]{Show of:}
\begin{itemize}
\item Instalando o meu ambiente numa VM de forma semi-automática
\begin{center}
\includegraphics[width=.9\linewidth]{./meme1.jpg}
\end{center}
\end{itemize}
\end{frame}

\section{Como usar Nix?}
\label{sec:orga8abfea}
\begin{frame}[label={sec:org5e4da66}]{Como usar Nix?}
\begin{itemize}
\item Basta instalar e usar!
\item Também podemos ter arquivos .nix para definirmos nossos shells.
\end{itemize}
\end{frame}

\section{Como usar Home Manager?}
\label{sec:org541084d}

\begin{frame}[label={sec:orgd46e7b7}]{Como usar Home Manager?}
\begin{itemize}
\item Basta editar os arquivos .nix
\item O que colocar aqui?
\begin{itemize}
\item Tudo que for relativo a algo específico de um usuário. Ex: vim.
\end{itemize}
\end{itemize}
\end{frame}

\section{Como usar NixOS?}
\label{sec:orgf8a0c0e}

\begin{frame}[label={sec:orgad9c757}]{Como usar NixOS?}
\begin{itemize}
\item Basta editar os arquivos .nix em /etc/nix!
\item Esse eu fizer :poop:? é só dar rollback!
\end{itemize}
\end{frame}

\section{Nix vs Docker}
\label{sec:org577b73e}
\begin{frame}[label={sec:org19e1682}]{Nix vs Docker}
\begin{itemize}
\item Nix não é um substituto para o Docker, ele faz coisas a mais e coisas a menos.
\item Coisas a mais: gerenciamento de dependências do sistema
\item Coisas a menos: não é um ambiente de container, não gerencia recurso do sistema\ldots{}
\end{itemize}
\end{frame}

\section{Vantagems e desvantagens do NixOS}
\label{sec:orga81393d}
\begin{frame}[label={sec:org8da6f29}]{Vantagems e desvantagens do NixOS}
\begin{itemize}
\item Vantagens:
\begin{enumerate}
\item Seu sistema é reproduzível em qualquer lugar
\item Declarativo
\item Modular
\item Facilita litar diferentes versões de pacotes
\item Versionamento
\end{enumerate}
\item Desvantagens:
\begin{enumerate}
\item É uma distro pouco usada, comunidade pequena, pouca documentação, mais problemas do que pessoas
\item Usa mais espaço
\item Usa bastante CPU para instalar pacotes
\item Curva de apredizado pesada
\end{enumerate}
\end{itemize}
\end{frame}

\section{Conclusão}
\label{sec:org1348505}
\begin{frame}[label={sec:orgffffbc7}]{Conclusão}
\begin{itemize}
\item IMHO é a melhor distro para desenvolvedores
\item Facilita a minha vida ao lidar com vários computadores
\item Considero esse ambiente mais confiável, pois eu posso dar rollback
\item A transição demora, mas é feita de forma suave
\item Como tudo é arquivo é mais fácil compartilhar coisas
\end{itemize}
\end{frame}
\end{document}
